\documentclass[11pt,letterpaper]{article}

\usepackage[utf8]{inputenc}
\usepackage[T1]{fontenc}
\usepackage[spanish,es-tabla]{babel}
\usepackage{geometry}
\usepackage{booktabs}
\usepackage{array}
\usepackage{longtable}
\usepackage{amsmath,amssymb}
\usepackage{xcolor}
\usepackage{enumitem}
\usepackage{hyperref}
\usepackage{fancyhdr}
\usepackage{lastpage}
\usepackage{titlesec}

\geometry{margin=2.2cm}
\hypersetup{
    colorlinks=true,
    linkcolor=black,
    urlcolor=blue
}

\definecolor{azul}{HTML}{1F4E79}
\definecolor{gris}{HTML}{F2F5F7}

\pagestyle{fancy}
\fancyhf{}
\lhead{Solemne IA Optimización 2026}
\rhead{Algoritmo Genético}
\cfoot{\thepage\ de \pageref{LastPage}}

\titleformat{\section}
  {\Large\bfseries\color{azul}}
  {\thesection.}
  {0.5em}
  {}

\titleformat{\subsection}
  {\large\bfseries}
  {\thesubsection}
  {0.5em}
  {}

\newcommand{\R}{\mathrm{R}}
\newcommand{\D}{\mathrm{D}}
\newcommand{\Depot}{0}
\newcommand{\DeltaMax}{0.30}

\begin{document}

\begin{titlepage}
    \centering
    \vspace*{1.2cm}
    {\Large Facultad de Ingeniería \par}
    \vspace{0.4cm}
    {\large CINF105 Optimización \par}
    \vspace{2.0cm}
    {\Huge\bfseries SOLEMNE IA \par}
    \vspace{0.4cm}
    {\LARGE Ruteo y scheduling de camiones cisterna \par}
    \vspace{0.4cm}
    {\Large Solución autosuficiente mediante Algoritmo Genético \par}
    \vfill
    \begin{tabular}{rl}
        Nombre: & \rule{8cm}{0.4pt} \\
        NRC: & \rule{8cm}{0.4pt} \\
        RUT: & \rule{8cm}{0.4pt}
    \end{tabular}
    \vfill
    {\large Junio 2026 \par}
\end{titlepage}

\tableofcontents
\newpage

\section*{Resumen metodológico}
\addcontentsline{toc}{section}{Resumen metodológico}

Este informe responde directamente los cuatro puntos de la evaluación y usa un único enfoque:
Algoritmo Genético (GA). La formulación matemática se presenta como el sistema de costos,
restricciones y penalizaciones que evalúa el GA. No se usan métodos comparativos ni archivos
externos para justificar la solución.

La idea central es codificar cada solución candidata como un cromosoma:
\[
\chi =
\left(
    \text{rutas por camión},
    \text{producto asignado a cada compartimento},
    \text{volúmenes cargados},
    \text{orden de carga en depósito}
\right).
\]
Cada cromosoma se evalúa con una función de costo total y penalizaciones por violar demanda,
capacidad, estabilidad, ventanas de tiempo o reglas de scheduling. Por lo tanto, el GA busca
cromosomas de menor penalización y menor costo real.

\section{Modelamiento determinista}

Se considera el escenario normal \(s=1\), con las demandas entregadas en el enunciado.
El depósito se denota como nodo \(0\), y las estaciones como \(J=\{1,2,3,4\}\).

\subsection{a) Conjuntos y parámetros}

\begin{longtable}{p{0.18\linewidth}p{0.75\linewidth}}
\toprule
\textbf{Símbolo} & \textbf{Descripción} \\
\midrule
\endhead
\(K=\{T1,T2\}\) & Conjunto de camiones. \\
\(C_k=\{C0,C1\}\) & Compartimentos del camión \(k\). \\
\(P=\{\R,\D\}\) & Productos: Regular y Diésel. \\
\(N=\{0,1,2,3,4\}\) & Nodos: depósito y estaciones. \\
\(J=N\setminus\{0\}\) & Estaciones de servicio. \\
\(d_{jp}\) & Demanda de la estación \(j\) para el producto \(p\). \\
\(D_{ij}\) & Distancia en kilómetros entre nodos \(i\) y \(j\). \\
\(\mathrm{cap}_{kc}\) & Capacidad del compartimento \(c\) del camión \(k\). \\
\(F_k\) & Costo fijo por usar el camión \(k\). \\
\([a_j,b_j]\) & Ventana de tiempo de la estación \(j\). \\
\(\tau_{ij}=D_{ij}/60\) & Tiempo de viaje en horas, usando 60 km/h. \\
\(h\) & Tiempo de servicio en estación: 30 minutos. \\
\(\Delta\) & Máxima diferencia permitida entre fill ratios: \(\Delta=0.30\). \\
\(\alpha\) & Costo por kilómetro: \(\alpha=2\). \\
\(\beta\) & Penalización por litro no entregado: \(\beta=10\). \\
\bottomrule
\end{longtable}

\subsection{b) Variables de decisión}

\begin{longtable}{p{0.21\linewidth}p{0.14\linewidth}p{0.58\linewidth}}
\toprule
\textbf{Variable} & \textbf{Tipo} & \textbf{Interpretación} \\
\midrule
\endhead
\(x_{kij}\) & Binaria & 1 si el camión \(k\) viaja directamente de \(i\) a \(j\). \\
\(r_k\) & Binaria & 1 si el camión \(k\) es utilizado. \\
\(z_{kj}\) & Binaria & 1 si el camión \(k\) visita la estación \(j\). \\
\(y_{kcp}\) & Binaria & 1 si el compartimento \(c\) del camión \(k\) transporta producto \(p\). \\
\(L_{kc}\) & Continua & Litros cargados inicialmente en el compartimento \(c\) del camión \(k\). \\
\(q_{kjp}\) & Continua & Litros del producto \(p\) entregados por el camión \(k\) en la estación \(j\). \\
\(u_{jp}\) & Continua & Shortage: litros no entregados del producto \(p\) en la estación \(j\). \\
\(t_{kj}\) & Continua & Hora de inicio de servicio del camión \(k\) en la estación \(j\). \\
\(\rho_{kcm}\) & Continua & Fill ratio del compartimento \(c\) del camión \(k\) después del evento \(m\) de su ruta. \\
\bottomrule
\end{longtable}

En la implementación GA, estas variables no se resuelven por simplex ni por branch-and-bound.
Se derivan desde el cromosoma \(\chi\), se reparan cuando es posible y se penalizan cuando
violan una restricción.

\subsection{c) Función objetivo}

El costo real de una solución factible es:
\[
    C(\chi)=
    \alpha \sum_{k\in K}\sum_{i\in N}\sum_{\substack{j\in N\\j\neq i}} D_{ij}x_{kij}
    + \sum_{k\in K}F_kr_k
    + \beta\sum_{j\in J}\sum_{p\in P}u_{jp}.
\]

Como el GA puede generar soluciones intermedias infactibles, se usa una función penalizada:
\[
    \Phi(\chi)=C(\chi)
    +\lambda_1V_{\text{flujo}}
    +\lambda_2V_{\text{demanda}}
    +\lambda_3V_{\text{capacidad}}
    +\lambda_4V_{\text{estabilidad}}
    +\lambda_5V_{\text{tiempo}}
    +\lambda_6V_{\text{scheduling}}.
\]

La aptitud se maximiza con:
\[
    \mathrm{fitness}(\chi)=\frac{1}{1+\Phi(\chi)}.
\]
Así, una solución con menor costo y menos violaciones tiene mayor probabilidad de ser
seleccionada para reproducción.

\subsection{d) Restricciones}

\subsubsection*{Conservación de flujo en routing}
Si un camión visita una estación, debe entrar y salir una vez:
\[
    \sum_{\substack{i\in N\\i\neq j}}x_{kij}=z_{kj},
    \qquad
    \sum_{\substack{i\in N\\i\neq j}}x_{kji}=z_{kj}
    \qquad \forall k\in K,\; j\in J.
\]
El uso del camión queda ligado a la salida y regreso al depósito:
\[
    \sum_{j\in J}x_{k0j}=r_k,
    \qquad
    \sum_{i\in J}x_{ki0}=r_k
    \qquad \forall k\in K.
\]

\subsubsection*{Satisfacción de demanda con shortage}
\[
    \sum_{k\in K}q_{kjp}+u_{jp}=d_{jp}
    \qquad \forall j\in J,\;p\in P,
\]
\[
    q_{kjp}\le d_{jp}z_{kj}
    \qquad \forall k\in K,\;j\in J,\;p\in P,
    \qquad
    u_{jp}\ge 0.
\]

\subsubsection*{Compatibilidad compartimento-producto}
Cada compartimento transporta a lo más un producto:
\[
    \sum_{p\in P}y_{kcp}\le 1
    \qquad \forall k\in K,\;c\in C_k.
\]
Solo se puede entregar un producto si existe un compartimento compatible en el camión:
\[
    q_{kjp}\le d_{jp}\sum_{c\in C_k}y_{kcp}
    \qquad \forall k\in K,\;j\in J,\;p\in P.
\]

\subsubsection*{Capacidad de compartimentos}
\[
    0\le L_{kc}\le \mathrm{cap}_{kc}\sum_{p\in P}y_{kcp}
    \qquad \forall k\in K,\;c\in C_k.
\]
Para cada compartimento asignado al producto \(p\):
\[
    \sum_{j\in J}q_{kjp}\le
    \sum_{c\in C_k}L_{kc}y_{kcp}
    \qquad \forall k\in K,\;p\in P.
\]

\subsubsection*{Estabilidad de carga}
Sea \(V_{kcm}\) el volumen restante del compartimento \(c\) del camión \(k\) después del
evento \(m\) de su ruta. El fill ratio es:
\[
    \rho_{kcm}=\frac{V_{kcm}}{\mathrm{cap}_{kc}}.
\]
La estabilidad exige:
\[
    |\rho_{kc m}-\rho_{kc' m}|\le \Delta
    \qquad
    \forall k\in K,\; c,c'\in C_k,\;m.
\]
En el GA, si esta desigualdad se viola en cualquier instante de la ruta, se suma la violación:
\[
    V_{\text{estabilidad}}=
    \sum_{k,m}\max\left\{0,\;|\rho_{kC0m}-\rho_{kC1m}|-\Delta\right\}.
\]

\subsubsection*{Ventanas de tiempo}
\[
    t_{ki}+h+\tau_{ij}
    \le t_{kj}+M(1-x_{kij})
    \qquad \forall k\in K,\;i,j\in J,\;i\neq j,
\]
\[
    a_jz_{kj}\le t_{kj}\le b_j+M(1-z_{kj})
    \qquad \forall k\in K,\;j\in J.
\]
Si el camión llega antes de \(a_j\), puede esperar; el servicio comienza en \(a_j\).

\section{Scheduling de carga en el depósito}

La carga es secuencial por camión y existe una bahía para Regular y una bahía para Diésel.
Los volúmenes fijos indicados en el enunciado son:

\begin{table}[h]
\centering
\begin{tabular}{llrrr}
\toprule
\textbf{Operación} & \textbf{Producto} & \textbf{Volumen (L)} & \textbf{Tasa (L/min)} & \textbf{Tiempo total (min)} \\
\midrule
T1-C0 & Regular & 5000 & 500 & \(5000/500+5=15.0\) \\
T1-C1 & Diésel & 5000 & 400 & \(5000/400+5=17.5\) \\
T2-C0 & Diésel & 4000 & 400 & \(4000/400+5=15.0\) \\
T2-C1 & Regular & 3500 & 500 & \(3500/500+5=12.0\) \\
\bottomrule
\end{tabular}
\caption{Cálculo de tiempos de carga incluyendo limpieza de 5 minutos.}
\end{table}

\subsection{a) Cálculo de tiempos de carga}

Los tiempos de proceso son:
\[
p_{T1,C0}=15,\quad
p_{T1,C1}=17.5,\quad
p_{T2,C0}=15,\quad
p_{T2,C1}=12.
\]

\subsection{b) Formulación del scheduling}

Sea \(I\) el conjunto de operaciones de carga. Para cada operación \(i\in I\), se define:
\[
    s_i=\text{inicio de la operación},\qquad
    p_i=\text{duración},\qquad
    B_i=\text{bahía requerida}.
\]
El objetivo es minimizar:
\[
    \min C_{\max}
\]
sujeto a:
\[
    s_i+p_i\le C_{\max}\qquad \forall i\in I.
\]
Precedencia dentro de cada camión:
\[
    s_{T1,C1}\ge s_{T1,C0}+15,
    \qquad
    s_{T2,C1}\ge s_{T2,C0}+15.
\]
No solapamiento en la misma bahía. Para dos operaciones \(i,j\) que usan la misma bahía:
\[
    s_i+p_i\le s_j+M(1-y_{ij}),
\]
\[
    s_j+p_j\le s_i+My_{ij},
\]
donde \(y_{ij}=1\) indica que \(i\) precede a \(j\).

\subsection{c) Secuencia óptima y diagrama de Gantt}

Una secuencia óptima es iniciar en paralelo T1-C0 en la bahía Regular y T2-C0 en la bahía
Diésel. Luego se cargan los segundos compartimentos de cada camión:

\begin{table}[h]
\centering
\begin{tabular}{lccc}
\toprule
\textbf{Bahía} & \textbf{Operación} & \textbf{Inicio relativo} & \textbf{Fin relativo} \\
\midrule
Regular & T1-C0 & 0.0 & 15.0 \\
Regular & T2-C1 & 15.0 & 27.0 \\
Diésel & T2-C0 & 0.0 & 15.0 \\
Diésel & T1-C1 & 15.0 & 32.5 \\
\bottomrule
\end{tabular}
\caption{Gantt esquemático de la secuencia óptima.}
\end{table}

El makespan mínimo es:
\[
    C_{\max}=32.5\text{ minutos}.
\]
Para que T1 salga a las 05:00, la carga debe comenzar a más tardar a las 04:27:30. Con ese
inicio, T1 queda listo exactamente a las 05:00 y T2 queda listo a las 04:54:30, por lo que
también puede salir a las 05:30.

\section{Análisis de la solución propuesta por el operador}

El operador propone:

\begin{table}[h]
\centering
\begin{tabular}{llllcc}
\toprule
\textbf{Camión} & \textbf{Ruta} & \textbf{C0} & \textbf{C1} & \textbf{Fill C0} & \textbf{Fill C1} \\
\midrule
T1 & \(D\to1\to3\to D\) & Regular & Diésel & 62.5\% & 71.4\% \\
T2 & \(D\to2\to4\to D\) & Diésel & Regular & 66.7\% & 38.9\% \\
\bottomrule
\end{tabular}
\end{table}

\subsection{a) Verificación de factibilidad}

\subsubsection*{Capacidad de compartimentos}

Los fill ratios declarados implican las siguientes cargas:
\[
T1:C0=0.625(8000)=5000\text{ L Regular},\quad
T1:C1=0.714(7000)=5000\text{ L Diésel}.
\]
\[
T2:C0=0.667(6000)=4000\text{ L Diésel},\quad
T2:C1=0.389(9000)=3500\text{ L Regular}.
\]

\begin{table}[h]
\centering
\begin{tabular}{lcccr}
\toprule
\textbf{Camión} & \textbf{Producto} & \textbf{Carga (L)} & \textbf{Demanda ruta (L)} & \textbf{Shortage mínimo} \\
\midrule
T1 & Regular & 5000 & \(3000+2500=5500\) & 500 \\
T1 & Diésel & 5000 & \(2000+3000=5000\) & 0 \\
T2 & Diésel & 4000 & \(1500+2500=4000\) & 0 \\
T2 & Regular & 3500 & \(4000+1000=5000\) & 1500 \\
\bottomrule
\end{tabular}
\caption{La solución no puede satisfacer toda la demanda: faltan al menos 2000 L.}
\end{table}

Por capacidad y carga declarada, la afirmación ``shortage = 0'' es falsa.

\subsubsection*{Estabilidad de carga}

Para T1, si se entrega lo posible en su ruta:

\begin{table}[h]
\centering
\begin{tabular}{lccc}
\toprule
\textbf{Instante} & \(\rho_{C0}\) Regular & \(\rho_{C1}\) Diésel & \textbf{Diferencia} \\
\midrule
Salida depósito & \(5000/8000=0.625\) & \(5000/7000=0.714\) & 0.089 \\
Después estación 1 & \(2000/8000=0.250\) & \(3000/7000=0.429\) & 0.179 \\
Después estación 3 & 0.000 & 0.000 & 0.000 \\
\bottomrule
\end{tabular}
\end{table}

T1 cumple estabilidad, pero no puede entregar 500 L de Regular en estación 3.

Para T2, si intenta entregar en la ruta \(D\to2\to4\to D\):

\begin{table}[h]
\centering
\begin{tabular}{lccc}
\toprule
\textbf{Instante} & \(\rho_{C0}\) Diésel & \(\rho_{C1}\) Regular & \textbf{Diferencia} \\
\midrule
Salida depósito & \(4000/6000=0.667\) & \(3500/9000=0.389\) & 0.278 \\
Después estación 2 & \(2500/6000=0.417\) & \(0/9000=0.000\) & 0.417 \\
Después estación 4 & 0.000 & 0.000 & 0.000 \\
\bottomrule
\end{tabular}
\end{table}

Después de estación 2 se obtiene \(0.417>0.30\), por lo que T2 viola la restricción de
estabilidad. Si se reduce la entrega de Regular en estación 2 para mantener estabilidad,
aumenta el shortage. En ambos casos, la solución propuesta es infactible.

\subsubsection*{Ventanas de tiempo}

Con velocidad de 60 km/h, cada kilómetro equivale a un minuto.

\begin{table}[h]
\centering
\begin{tabular}{llcl}
\toprule
\textbf{Camión} & \textbf{Tramo} & \textbf{Llegada sin espera} & \textbf{Ventana} \\
\midrule
T1 & Depósito \(\to\) 1 & 05:20 & [06:00, 10:00] \\
T1 & 1 \(\to\) 3 & 06:15 & [08:00, 14:00] \\
T2 & Depósito \(\to\) 2 & 06:05 & [07:00, 12:00] \\
T2 & 2 \(\to\) 4 & 06:50 & [06:00, 11:00] \\
\bottomrule
\end{tabular}
\end{table}

Las llegadas crudas a 1, 3 y 2 son tempranas. Si el modelo permite espera, el servicio puede
iniciarse dentro de las ventanas; si se exige llegada estrictamente dentro de la ventana, también
hay violación temporal. En cualquier interpretación, la solución ya es infactible por capacidad
y estabilidad.

\subsection{b) Verificación del costo declarado}

La distancia real recorrida es:
\[
T1: 20+25+15=60\text{ km},\qquad
T2: 35+15+40=90\text{ km}.
\]
\[
\text{Distancia total}=150\text{ km},\qquad
\text{Costo distancia}=2(150)=300.
\]

\begin{table}[h]
\centering
\begin{tabular}{lrr}
\toprule
\textbf{Componente} & \textbf{Cálculo} & \textbf{Costo} \\
\midrule
Distancia & \(150\times2\) & 300 \\
Costos fijos & \(500+400\) & 900 \\
Shortage mínimo por carga declarada & \(2000\times10\) & 20000 \\
\midrule
\textbf{Total real mínimo} & & \textbf{21200} \\
\bottomrule
\end{tabular}
\end{table}

El costo declarado de 260 es incorrecto: omite costos fijos, subestima la distancia y declara
shortage cero aunque faltan al menos 2000 litros.

\subsection{c) Mejora inicial y solución final encontrada por GA}

Como primer acercamiento, una mejora simple es mantener T1 en \(D\to1\to3\to D\),
corregir las cargas reales requeridas y cambiar el orden de T2 a \(D\to4\to2\to D\).
Esta modificación elimina el shortage y corrige la estabilidad:

\begin{table}[h]
\centering
\begin{tabular}{llllcc}
\toprule
\textbf{Camión} & \textbf{Ruta} & \textbf{C0} & \textbf{C1} & \textbf{Carga C0} & \textbf{Carga C1} \\
\midrule
T1 & \(D\to1\to3\to D\) & Regular & Diésel & 5500 & 5000 \\
T2 & \(D\to4\to2\to D\) & Diésel & Regular & 4000 & 5000 \\
\bottomrule
\end{tabular}
\end{table}

Verificación de estabilidad:

\begin{table}[h]
\centering
\begin{tabular}{lccc}
\toprule
\textbf{Camión e instante} & \(\rho_{C0}\) & \(\rho_{C1}\) & \textbf{Diferencia} \\
\midrule
T1 salida & \(5500/8000=0.688\) & \(5000/7000=0.714\) & 0.027 \\
T1 después estación 1 & \(2500/8000=0.313\) & \(3000/7000=0.429\) & 0.116 \\
T1 después estación 3 & 0.000 & 0.000 & 0.000 \\
T2 salida & \(4000/6000=0.667\) & \(5000/9000=0.556\) & 0.111 \\
T2 después estación 4 & \(1500/6000=0.250\) & \(4000/9000=0.444\) & 0.194 \\
T2 después estación 2 & 0.000 & 0.000 & 0.000 \\
\bottomrule
\end{tabular}
\end{table}

Todas las diferencias son menores o iguales a \(\Delta=0.30\).

Verificación de ventanas con espera:

\begin{table}[h]
\centering
\begin{tabular}{llll}
\toprule
\textbf{Camión} & \textbf{Estación} & \textbf{Llegada / servicio} & \textbf{Cumple} \\
\midrule
T1 & 1 & llega 05:20, espera, sirve 06:00--06:30 & Sí \\
T1 & 3 & llega 06:55, espera, sirve 08:00--08:30 & Sí \\
T2 & 4 & llega 06:10, sirve 06:10--06:40 & Sí \\
T2 & 2 & llega 06:55, espera, sirve 07:00--07:30 & Sí \\
\bottomrule
\end{tabular}
\end{table}

El costo de esta mejora inicial es:
\[
\text{Distancia}=60+(40+15+35)=150\text{ km},
\]
\[
\text{Costo total}=2(150)+500+400+10(0)=1200.
\]
Por lo tanto, la mejora inicial reduce el costo real mínimo desde 21200 a 1200. Sin
embargo, al ejecutar el Algoritmo Genético sobre cromosomas de rutas, asignación de
productos y cargas reparadas, se encuentra una solución factible de menor costo:

\begin{table}[h]
\centering
\begin{tabular}{llllcc}
\toprule
\textbf{Camión} & \textbf{Ruta GA} & \textbf{C0} & \textbf{C1} & \textbf{Carga C0} & \textbf{Carga C1} \\
\midrule
T1 & \(D\to1\to2\to4\to D\) & Regular & Diésel & 8000 & 6000 \\
T2 & \(D\to3\to D\) & Diésel & Regular & 3000 & 2500 \\
\bottomrule
\end{tabular}
\caption{Solución final encontrada por el Algoritmo Genético.}
\end{table}

La verificación de estabilidad de la solución GA es:

\begin{table}[h]
\centering
\begin{tabular}{lccc}
\toprule
\textbf{Camión e instante} & \(\rho_{C0}\) & \(\rho_{C1}\) & \textbf{Diferencia} \\
\midrule
T1 salida & \(8000/8000=1.000\) & \(6000/7000=0.857\) & 0.143 \\
T1 después estación 1 & \(5000/8000=0.625\) & \(4000/7000=0.571\) & 0.054 \\
T1 después estación 2 & \(1000/8000=0.125\) & \(2500/7000=0.357\) & 0.232 \\
T1 después estación 4 & 0.000 & 0.000 & 0.000 \\
T2 salida & \(3000/6000=0.500\) & \(2500/9000=0.278\) & 0.222 \\
T2 después estación 3 & 0.000 & 0.000 & 0.000 \\
\bottomrule
\end{tabular}
\end{table}

Todas las diferencias son menores o iguales a \(\Delta=0.30\). Además, no hay shortage:
T1 entrega Regular y Diésel a las estaciones 1, 2 y 4, mientras T2 entrega Diésel y
Regular a la estación 3.

\begin{table}[h]
\centering
\begin{tabular}{llll}
\toprule
\textbf{Camión} & \textbf{Estación} & \textbf{Llegada / servicio} & \textbf{Cumple} \\
\midrule
T1 & 1 & llega 05:20, espera, sirve 06:00--06:30 & Sí \\
T1 & 2 & llega 06:40, espera, sirve 07:00--07:30 & Sí \\
T1 & 4 & llega 07:45, sirve 07:45--08:15 & Sí \\
T2 & 3 & llega 05:45, espera, sirve 08:00--08:30 & Sí \\
\bottomrule
\end{tabular}
\caption{Ventanas de tiempo de la solución GA.}
\end{table}

El costo final de la solución GA es:
\[
\text{Distancia}=20+10+15+40+15+15=115\text{ km},
\]
\[
\text{Costo total}=2(115)+500+400+10(0)=1130.
\]
Esta es la solución adoptada como resultado final del informe: mejora el primer acercamiento
de 1200, mantiene factibilidad de capacidad, estabilidad y ventanas de tiempo, y elimina el
shortage.

\section{Extensión estocástica}

La demanda real se modela mediante tres escenarios. Cuando el enunciado no modifica el Diésel
de estación 3, se conserva el valor normal de 3000 L.

\begin{table}[h]
\centering
\begin{tabular}{lrrrrrrrr}
\toprule
\textbf{Esc.} & \textbf{Prob.} & \textbf{E1 R} & \textbf{E1 D} & \textbf{E2 R} & \textbf{E2 D} & \textbf{E3 R} & \textbf{E3 D} & \textbf{E4 R/D} \\
\midrule
s=1 & 0.5 & 3000 & 2000 & 4000 & 1500 & 2500 & 3000 & 1000 / 2500 \\
s=2 & 0.3 & 4500 & 2800 & 5500 & 2000 & 3500 & 3000 & 1000 / 2500 \\
s=3 & 0.2 & 2000 & 1200 & 3000 & 1000 & 1800 & 3000 & 1000 / 2500 \\
\bottomrule
\end{tabular}
\end{table}

\subsection{a) Modelo de dos etapas}

\subsubsection*{Decisiones de primera etapa}
Son las decisiones tomadas antes de conocer la demanda real:
\begin{itemize}[leftmargin=*]
    \item Activar o no cada camión: \(r_k\).
    \item Asignar producto a cada compartimento: \(y_{kcp}\).
    \item Definir la estructura de ruta o prioridad de visita de cada camión: \(x_{kij}\).
    \item Definir el scheduling de carga en depósito: inicios \(s_i\) y orden de bahía.
\end{itemize}

\subsubsection*{Decisiones de segunda etapa}
Después de observar el escenario \(s\), se ajustan:
\begin{itemize}[leftmargin=*]
    \item Volúmenes efectivamente entregados \(q_{kjp}^{s}\).
    \item Shortage por estación y producto \(u_{jp}^{s}\).
    \item Tiempos de servicio y espera \(t_{kj}^{s}\).
    \item Fill ratios y verificación de estabilidad \(\rho_{kcm}^{s}\).
\end{itemize}

\subsubsection*{Función objetivo estocástica}
Para un cromosoma de primera etapa \(\chi\), el GA evalúa el costo esperado:
\[
    \min_{\chi}\;
    C^{(1)}(\chi)
    +\sum_{s\in S}\pi_s Q_s(\chi),
\]
donde:
\[
    C^{(1)}(\chi)=
    \sum_{k\in K}F_kr_k
    +V_{\text{scheduling}}(\chi),
\]
\[
    Q_s(\chi)=
    \alpha\sum_{k,i,j}D_{ij}x_{kij}
    +\beta\sum_{j,p}u_{jp}^{s}
    +\lambda_1V_{\text{demanda}}^{s}
    +\lambda_2V_{\text{capacidad}}^{s}
    +\lambda_3V_{\text{estabilidad}}^{s}
    +\lambda_4V_{\text{tiempo}}^{s}.
\]

\subsubsection*{Restricciones de primera y segunda etapa}
Primera etapa:
\[
    r_k,\; y_{kcp},\; x_{kij},\; s_i
    \quad \text{no dependen de }s.
\]
Esto es la condición de no anticipatividad: no se puede cambiar la ruta base o la asignación
de compartimentos antes de conocer el escenario.

Segunda etapa, para cada escenario \(s\):
\[
    \sum_{k\in K}q_{kjp}^{s}+u_{jp}^{s}=d_{jp}^{s}
    \qquad \forall j,p,s,
\]
\[
    \sum_{j\in J}q_{kjp}^{s}\le
    \sum_{c\in C_k}L_{kc}^{s}y_{kcp}
    \qquad \forall k,p,s,
\]
\[
    |\rho_{kC0m}^{s}-\rho_{kC1m}^{s}|\le 0.30
    \qquad \forall k,m,s,
\]
\[
    a_jz_{kj}\le t_{kj}^{s}\le b_j+M(1-z_{kj})
    \qquad \forall k,j,s.
\]

En términos de GA, cada cromosoma se prueba contra los tres escenarios y su aptitud final
es el costo esperado penalizado. Así se favorecen soluciones robustas: no solo baratas en el
escenario normal, sino también menos costosas cuando la demanda sube o baja.

\section*{Conclusión}
\addcontentsline{toc}{section}{Conclusión}

La solución del operador no es factible: las cargas declaradas no alcanzan para satisfacer la
demanda, T2 viola la estabilidad después de estación 2 y el costo declarado omite componentes
obligatorios. La ruta \(D\to4\to2\to D\) para T2 sirve como primer acercamiento factible
con costo 1200, pero el Algoritmo Genético encuentra una solución mejor: T1 atiende
\(1\to2\to4\), T2 atiende \(3\), no existe shortage, se respetan las ventanas de tiempo y
el costo total baja a 1130.

El Algoritmo Genético permite abordar el problema completo codificando rutas, asignación de
productos, volúmenes y scheduling en un cromosoma único. La función de aptitud combina costo
real con penalizaciones, por lo que el método puede evaluar tanto el escenario determinista como
la extensión estocástica de dos etapas sin recurrir a métodos comparativos.

\end{document}
